% 图6(a) 车道汇入·Baseline：汇入车 ST 轨迹为倾斜平行四边形，
% 全时刻并集投影把其左沿保守压至 t=0，自车从起点即被挡，被迫停车。
\begin{tikzpicture}[scale=1.0, transform shape, arr/.style={-Stealth, thick}]
  \useasboundingbox (-0.75,-1.05) rectangle (5.85,4.75);

  % 汇入车 ST 轨迹：倾斜平行四边形，左沿压至 t=0（左下角 (0,1.2)）
  \fill[dangerred!35, draw=dangerred, thick]
    (0,1.2) -- (4.0,3.3) -- (4.0,4.3) -- (0,2.2) -- cycle;
  \node[font=\scriptsize\bfseries, dangerred!85!black, rotate=27] at (2.0,3.0) {汇入车 ST 轨迹};
  % 压至 t=0 标注
  \draw[dangerred, very thick] (0,1.2) -- (0,2.2);
  \node[font=\tiny, dangerred!80!black, anchor=east, align=right] at (-0.05,1.7) {左沿\\压至\\$t{=}0$};

  % 坐标轴
  \draw[arr] (0,0) -- (5.6,0) node[right, font=\scriptsize] {$t$/s};
  \draw[arr] (0,0) -- (0,4.5) node[above, font=\scriptsize] {$s$/m};
  \node[font=\tiny, left=2pt] at (0,0) {$0$};

  % Baseline s(t)：起步即撞平行四边形下沿，被迫停车（几乎不前进）
  \draw[deepblue, very thick, -Stealth]
    (0,0) .. controls (0.3,0.6) and (0.6,0.95) .. (0.8,1.0) -- (5.0,1.0);
  \fill[deepblue] (0.8,1.0) circle (1.6pt);
  \node[font=\tiny, deepblue!80!black, anchor=north] at (2.8,0.97) {自车被迫停车（无通行区）};
\end{tikzpicture}
